% !TEX root = ../Main.tex
\chapter{例题}

\section{菱形侧面 $\perp$ 底面}

\ex{
    如图，$D$ 是 $AC$ 的中点，四边形 $BDEF$ 是菱形，平面 $BDEF\perp$ 平面 $ABC$，$\angle FBD=60^\circ$，$AB\perp BC$，$AB=BC=\sqrt2$。
    \begin{enumerate}
        \item[(1)] 若点 $M$ 是线段 $BF$ 的中点，证明：$BF\perp$ 平面 $AMC$；
        \item[(2)] 求平面 $AEF$ 与平面 $BCF$ 所成的锐二面角的余弦值。
    \end{enumerate}
}

\section{侧面 $\perp$ 底面的三棱柱}

\ex{
    如图，已知三棱柱 $ABC$-$A_1B_1C_1$，侧面 $BCC_1B_1\perp$ 底面 $ABC$。
    \begin{enumerate}
        \item[(1)] 若 $M,N$ 分别是 $AB,A_1C$ 的中点，求证：$MN\parallel$ 平面 $BCC_1B_1$；
        \item[(2)] 若三棱柱 $ABC$-$A_1B_1C_1$ 的各棱长均为 $2$，侧棱 $BB_1$ 与底面 $ABC$ 所成的角为 $60^\circ$，问在线段 $A_1C_1$ 上是否存在一点 $P$，使得平面 $B_1CP\perp$ 平面 $ACC_1A_1$？若存在，求 $C_1P$ 与 $PA_1$ 的比值；若不存在，说明理由。
    \end{enumerate}
}

\section{矩形折叠成直二面角}

\ex{
    如图，在矩形 $ABCD$ 中，$AB=3$，$BC=3\sqrt3$，点 $E,H$ 分别是所在边靠近 $B,D$ 的三等分点，$O$ 是 $EH$ 的中点。现沿着 $EH$ 将矩形折成直二面角，分别连接 $AD,AC,CB$ 形成如图所示的多面体。
    \begin{enumerate}
        \item[(1)] 证明：$EH\perp AC$；
        \item[(2)] 求二面角 $B$-$AC$-$O$ 的平面角的余弦值。
    \end{enumerate}
}

\section{正三棱柱：题 20}

\ex{
    如图，三棱柱 $ABC$-$A_1B_1C_1$ 的底面是正三角形，侧面 $BB_1C_1C$ 是矩形，$M,N$ 分别为 $BC,B_1C_1$ 的中点，$P$ 为 $AM$ 上一点，过 $B_1C_1$ 和 $P$ 的平面交 $AB$ 于 $E$、交 $AC$ 于 $F$。
    \begin{enumerate}
        \item[(1)] 证明：$AA_1\parallel MN$，且平面 $A_1AMN\perp$ 平面 $EB_1C_1F$；
        \item[(2)] 设 $O$ 为 $\triangle A_1B_1C_1$ 的中心，若 $AO\parallel$ 平面 $EB_1C_1F$ 且 $AO=AB$，求直线 $B_1E$ 与平面 $A_1AMN$ 所成角的正弦值。
    \end{enumerate}
}

\section{菱形底 + 竖立矩形}

\ex{
    如图，菱形 $ABCD$ 的边长为 $4$，$\angle DAB=60^\circ$，矩形 $BDFE$ 的面积为 $8$，且平面 $BDFE\perp$ 平面 $ABCD$。
    \begin{enumerate}
        \item[(1)] 证明：$AC\perp BE$；
        \item[(2)] 求二面角 $E$-$AF$-$D$ 的正弦值。
    \end{enumerate}
}

\section{直三棱柱中的垂直与线面角}

\ex{
    如图，直三棱柱 $ABC$-$A_1B_1C_1$ 中，$AC=BC$，$AA_1=2$，$AB=\sqrt2$，$D$ 为 $BB_1$ 的中点，$E$ 为线段 $AB_1$ 上一点。
    \begin{enumerate}
        \item[(1)] 若 $DE\perp CD$，求证：$DE\perp AB_1$；
        \item[(2)] 若 $AE=2EB_1$，异面直线 $AB_1$ 与 $CD$ 所成的角为 $30^\circ$，求直线 $DE$ 与平面 $AA_1C_1C$ 所成角的正弦值。
    \end{enumerate}
}

\section{矩形 $ABCD$ + 垂面 $BEC$}

\ex{
    如图，在几何体 $ABCDE$ 中，四边形 $ABCD$ 是矩形，$AB\perp$ 平面 $BEC$，$BE\perp EC$，$AB=BE=EC=2$，$G$ 是线段 $BE$ 的中点，点 $F$ 在线段 $CD$ 上且 $GF\parallel$ 平面 $ADE$。
    \begin{enumerate}
        \item[(1)] 求 $CF$ 的长；
        \item[(2)] 求平面 $AEF$ 与平面 $AFG$ 的夹角的余弦值。
    \end{enumerate}
}

\section{四棱锥：等腰侧棱定垂面}

\ex{
    （12 分）如图，四棱锥 $P$-$ABCD$ 的底面 $ABCD$ 为矩形，$E$ 是边 $AD$ 上的点，$PA=AB=AE=2DE$，$\angle PBA=\angle PBC=60^\circ$。
    \begin{enumerate}
        \item[(1)] 求证：平面 $PBE\perp$ 平面 $ABCD$；
        \item[(2)] 求直线 $PC$ 和平面 $PBD$ 所成角的正弦值。
    \end{enumerate}
}

\section{四棱锥：直角梯形底}

\ex{
    如图所示，四棱锥 $P$-$ABCD$ 中，$PA\perp$ 平面 $ABCD$，$AD\parallel BC$，$AD\perp CD$，且 $AD=CD=2\sqrt2$，$BC=4\sqrt2$，$PA=2$，点 $M$ 在线段 $PD$ 上。
    \begin{enumerate}
        \item[(1)] 求证：$AB\perp PC$；
        \item[(2)] 若二面角 $M$-$AC$-$D$ 的大小为 $45^\circ$，求 $BM$ 与平面 $PAC$ 所成的角的正弦值。
    \end{enumerate}
}

\section{矩形沿 $AE$ 折叠}

\ex{
    如图 1，矩形 $ABCD$ 中，$AB=3\sqrt5$，$BC=2\sqrt5$，点 $E$ 在线段 $DC$ 上且 $DE=\sqrt5$。现将 $\triangle AED$ 沿 $AE$ 折到 $\triangle AED'$ 的位置，连接 $CD',BD'$（如图 2）。
    \begin{enumerate}
        \item[(1)] 若点 $P$ 在线段 $BC$ 上，且 $BP=\dfrac{\sqrt5}2$，证明：$AE\perp D'P$；
        \item[(2)] 记平面 $AD'E$ 与平面 $BCD'$ 的交线为 $l$。若二面角 $B$-$AE$-$D'$ 为 $\dfrac{2\pi}3$，求 $l$ 与平面 $D'CE$ 所成角的正弦值。
    \end{enumerate}
}
